% Included at the end of Section 6, after the complete observable construction.
\subsection{Adjoint reduction of the second-order-field propagation products}\label{sec:prop-adjoint}
The part of the second-order ray driven by $H_2$ can be eliminated exactly. This subsection isolates \emph{all terms linear in the second-order fields}
$H_2$ and $\nu_2$, excluding the top-order fields $H_3,\nu_3$ and the remaining products made
only of first-order fields. It is a decomposition of the already complete map, not an omission
or a new initial-state assumption. The adjoint below holds at general momenta and finite times.

Write $h=H_1$, $j=H_2$, $m=\nu_2$, $P=1-nn^T$, $e=\bar t_\star$, $o=\bar t_o$,
$\bar X=n(o-t)$. In this subsection $h',j'$ denote partial time derivatives before pullback;
derivatives of ray and adjoint variables are total derivatives along $\bar X$.
Spatial derivatives act only on the fields. Let $X=X_1$, $v=v_1$ be the known first-order
orthonormal ray, and let $X_j,v_j$ denote only the part of $X_2,v_2$ linear in $j$.
Since $\zeta^{(2)}=0$, their equations and observer data are
\begin{equation}\label{eq:prop-linear-ray}
 X_j'=v_j+\tfrac12jn,\qquad
 v_j'=G_j=\tfrac12Pj'n+\tfrac12[\nabla j_{nn}-(n\cdot\partial)(jn)],\qquad
 X_j(o)=v_j(o)=0.
\end{equation}
The lapse $m$ does not deflect conformal null rays; it enters this sector only through the
physical emission surface and observer clock. Put
\begin{equation}\label{eq:prop-bilinear}
 \mathfrak B[h,j]=\int_e^o\left[
 \tfrac12X\cdot\nabla j'_{nn}-j'(n,v)
 +\tfrac12X_j\cdot\nabla h'_{nn}-h'(n,v_j)\right]\dd t.
\end{equation}
The cubic bulk contribution linear in $j$ is $3\mathfrak B[h,j]$ in the full-field Taylor
convention. Kernel polarisation instead gives the sum over its three single--pair partitions,
with no extra factor three.

\begin{proposition}[Propagation adjoint]\label{prop:propagation-adjoint}
Define first-order-only functions
\begin{gather*}
 r=\tfrac12\nabla h'_{nn},\qquad s=-Ph'n,\qquad \boldsymbol Y=v+\tfrac12P\boldsymbol K,\\
 \boldsymbol R(t)=\int_e^t r(u)\dd u,\qquad
 \boldsymbol K(t)=\int_e^t[s(u)-\boldsymbol R(u)]\dd u.
\end{gather*}
Then the unknown second-order ray can be eliminated exactly:
\begin{equation}\label{eq:prop-adjoint-first}
 \mathfrak B[h,j]=\int_e^o\left[
 \tfrac12X\cdot\nabla j'_{nn}-j'(n,v)
 -\tfrac12\boldsymbol R\cdot jn-\boldsymbol K\cdot G_j\right]\dd t.
\end{equation}
A further integration by parts removes all partial time derivatives of $j$ from the bulk:
\begin{equation}\label{eq:prop-adjoint-bulk}
\boxed{\begin{aligned}
 \mathfrak B[h,j]={}&
 [\tfrac12X\cdot\nabla j_{nn}-\boldsymbol Y\cdot jn]_e^o\\
 &+\int_e^o\left\{
 \tfrac12X\cdot\nabla(n\cdot\partial)j_{nn}
 -\tfrac12(X'+P\boldsymbol K)\cdot\nabla j_{nn}\right.\\
 &\hspace{19mm}\left.-v\cdot(n\cdot\partial)(jn)
 +(\boldsymbol Y'-\tfrac12\boldsymbol R)\cdot jn\right\}\dd t.
\end{aligned}}
\end{equation}
Only first-order ray and adjoint functions and the second-order metric remain. In particular,
$\boldsymbol R$ must not be projected transversely in its coupling to $jn$: the coordinate
displacement has a longitudinal part even though $v_j$ is transverse.
\end{proposition}
\begin{proof}
Differentiate $\boldsymbol R\cdot X_j+\boldsymbol K\cdot v_j$ and use
\eqref{eq:prop-linear-ray}. Its derivative is
$r\cdot X_j+s\cdot v_j+\boldsymbol R\cdot jn/2+\boldsymbol K\cdot G_j$.
Its endpoint values vanish, because the adjoints vanish at $e$ and the ray responses at $o$.
This proves \eqref{eq:prop-adjoint-first}. In the remaining $j'$ terms use
$\partial_tj=\dd j/\dd t+(n\cdot\partial)j$ on the background ray. The longitudinal
$\boldsymbol K$ pieces cancel between the two spatial-derivative terms, leaving
\eqref{eq:prop-adjoint-bulk}. The resulting metric boundary is \emph{not} zero:
at $o$, $X=v=0$ but $\boldsymbol Y=P\boldsymbol K/2$ need not vanish.
\end{proof}

\paragraph{Emission, physical clock and normalisation.}
Define $F_h=h'_{nn}/2$, $F_j=j'_{nn}/2$,
$U_h=\int_e^oF_h\dd t$, $U_j=\int_e^oF_j\dd t$ and
\[
 \mathcal C[m]=\frac{-\bar t_I^2m(\bar t_I,0)-\int_{\bar t_I}^otm(t,0)\dd t}{o}.
\]
This is the part of $D_2$ linear in $m$, not an additional clock correction.
The complete composed-exponent contribution of this sector is
\begin{equation}\label{eq:prop-complete-sector}
 \widehat I_3\big|_{2}=3\mathfrak D_2,\qquad
 \mathfrak D_2=\mathfrak B[h,j]
 +[eF_hm-F_jd_1]_e+D_1\partial_oU_j+\mathcal C[m]\partial_oU_h.
\end{equation}
Here $d_1=-e\nu_1(e,\bar X_e)$, and $D_1$ is the original physical-clock coefficient.
The $\partial_o$ derivatives include all background-ray and upper-limit dependence, with
$e$ fixed. After centring, the corresponding full normalised-map contribution is
\begin{equation}\label{eq:prop-normalized-sector}
 \Theta^{(3)}\big|_2=3\mathsf P_0\left[
 -\mathfrak D_2+(U_h-\sky{U_h})(U_j-\sky{U_j})\right].
\end{equation}
Thus the second-order-map normalisation product cannot be discarded. The dipole is still
retained in that product. All terms made only of first-order fields, including the mixed
clock--emission product in \eqref{eq:cubic-composed}, remain outside this isolated sector and
must be added as before. The Fermi--Walker correction is also first-order-field-only at this
map order; it does not add a $j$-dependent term here.

\paragraph{An evaluated shell emission contraction.}
The emission part of this sector already fits the nine-moment construction. Let
$\mathcal C_S(t)=\langle q_SF_S+q_TF_T\rangle_\mu$ (so $\Upsilon=\mathsf a\mathcal C_S$),
$Q_k=\hat k\hat k-\delta/3$. After polarising the full cubic expression and averaging over
$\hat q$, the emission-exponent kernel is exactly
\begin{equation}\label{eq:prop-emission-moments}
 \overline{\mathcal I_{3,e}\big|_2}(k;n)
 =\frac e2\left[
 \mathsf b'(n_h+3N_2-N_0)+3\mathsf a\mathcal C_S'
 \right]_e Q_{k,nn}\,e^{ik\cdot n(o-e)}.
\end{equation}
Indeed the two mixed partitions share one emission time, so their phases combine to the
external phase. The azimuthal averages are
$\overline{Q_q}=P_2(\mu)Q_k$ and
$\overline{H_2(k,q)}=\tfrac32(q_SF_S+q_TF_T)Q_k$.
The homogeneous tensor averages to zero and the homogeneous lapse remains.
For $\ell\ge2$, its temperature-transfer contribution is
\[
 \mathcal A_{\ell,e}\big|_2
 =\frac e2 w_\ell(k(o-e))
 [\mathsf b'(n_h+3N_2-N_0)+3\mathsf a\mathcal C_S']_e.
\]
The sign follows from $\Theta=-I$ at linear order in an exponent kernel and the
$Q_{k,nn}$ Funk--Hecke projection. This contribution can be added to the explicit emission
boundary of \eqref{eq:adjoint-nine-moment}; no identity setting their sum to zero is asserted.

\paragraph{What this does and does not settle.}
The bulk reduction and complete sector \eqref{eq:prop-complete-sector} are checked exactly
in \path{solution/check_propagation_adjoint.py}: matrix-exponential driver, arbitrary
plane-wave jets, ray duality, boundary and clock coefficients, centring algebra and shell
emission projections, with omission controls. It does not evaluate the full propagated
$13$ response. In a mixed-time product, for example,
\[
 (k+q)\cdot n(o-t)-q\cdot n(o-u)
 =k\cdot n(o-t)+q\cdot n(u-t).
\]
The second phase is present unless $u=t$. Consequently the finite nine polynomial moments
proved for the top-order fields and shared-time emission terms cannot simply be assumed to
close the remaining propagation integrals. No cancellation of the full cubic response or
late-time limit is proved. The unevaluated remainder is the combination of
\eqref{eq:adjoint-nine-moment}, \eqref{eq:prop-adjoint-bulk}, the clock terms and the remaining
first-order-only products, with the displayed physical boundaries retained.
